\chapter*{Conclusions and Future Work}
\label{cap:conclusions}
\addcontentsline{toc}{chapter}{Conclusions and Future Work}
\markboth{Conclusions and Future Work}{Conclusions and Future Work}

\chapterquote{What we call the beginning is often the end. The end is where we start from.}{T. S. Eliot}

This chapter closes the memoir by gathering what the work has contributed, explicitly
acknowledging its limits, and stating the natural continuation lines that open from it. The
first section synthesizes the results achieved with respect to the objectives stated in the
Introduction. The second section delimits the scope of the delivered system. The third section
groups, in several complementary directions, the extensions left open as project continuations.

\section*{Conclusions}

The work has made it possible to build a digital collection exploration engine that integrates
the three mechanisms planned from the beginning: faceted filtering over metadata, transversal
navigation between collections through semantic references, and set operations over partial
results. Each one has materialized into an identifiable component of the system (filters,
\textit{links}, and the union set) and its behavior has been verified over real dumps of two
public catalogs of very different nature, TMDB (audiovisual) and DBLP (scientific publications),
chosen precisely for their structural differences.

The two-phase methodological strategy proved particularly sound. The first phase, centered on
the Java prototype without a graphical interface, allowed us to validate the technical viability
of the navigation model rigorously, without contaminating design decisions with constraints
derived from presentation. The second phase, devoted to exposing the engine through a
\textit{RESTful API} and developing the web interface, could then concentrate on user experience
matters on top of an already stable base. This clear separation between back-end and front-end
also facilitated pair work, by defining explicit contracts between the two layers.

The adoption of the \textbf{Scrum} framework adapted to a small team has shown its value by
allowing us to reorient the project whenever the findings of each \textit{sprint} forced the
reconsideration of previous decisions. In parallel, the usability evaluation conducted with
external participants following a formal protocol has produced qualitative evidence that has
guided the last interface iteration and constitutes, in itself, a reusable contribution for
future developments on the same engine.

Beyond the concrete components of the system, we consider that the most reusable contribution
of the work is the \textbf{formal model} underlying the engine: the precise definition of the
exploration state, the operations that transform it, and the invariants that guarantee its
coherence. This model responds to the gap identified in the State of the Art chapter, where
we observed that existing tools offer the three mechanisms (faceted filtering, navigation
between collections, and set operations) in isolation and without a unifying framework that
articulates them as pieces of the same calculus. Since the model has remained decoupled from
the implementation language and from the specific catalog, it can be reused as the foundation
of future engines built on other technologies or applied to other domains, regardless of the
particular decisions we have made in the delivered prototype.

\section*{Limitations}

The scope of the delivered system must be understood within several well-defined limits. The
functional validation of the engine was carried out over two distinct domains (the audiovisual
catalog TMDB and the bibliographic database DBLP), which provides reasonable evidence that the
model does not depend on the particularities of a single corpus, but does not exhaust the
spectrum of collections to which the system aspires to apply. Its extension to other types of
holdings (institutional document repositories, complete library catalogs, or open-source
software collections) remains, therefore, a reasonable hypothesis not empirically demonstrated
in this work. Similarly, the load tests have been limited to collections of a few tens of
thousands of entities, without addressing large-scale scenarios that would require specific
optimization work.

The system lacks, by design decision and not by oversight, a persistent user layer: sessions
are ephemeral, and identification, authentication, and state synchronization across devices are
not contemplated. The web interface, in turn, has been optimized for desktop displays and,
although it collapses reasonably on narrow screens, the mobile case has not been worked on
specifically. Finally, the usability evaluation follows a qualitative approach with a small
sample (five participants), suitable for detecting most observable problems
but insufficient to sustain quantitative claims of statistical character.

\section*{Future work}

From the state in which the system is delivered, we identify several natural continuation
lines. The first two (language-model-based assistance and engine distribution) have an applied
research component and would open the door to a larger-scale project. The third (incremental
query generation) affects the core of the navigation model and raises an interesting trade-off
between conceptual simplicity and execution cost. The remaining ones (user layer, extension
to new domains, multiple selection, reinforcement of the test suite, and user experience
improvements) are more incremental in nature and could be addressed independently, without
altering the navigation model's core.

\subsection*{Language-model-based assistance}

The emergence of large language models opens an interesting avenue to enrich the engine without
replacing it. Rather than delegating navigation to an opaque model, the model can be used as an
assistance layer: translation of natural-language queries into concrete filter combinations,
suggestion of pertinent filters depending on the exploration state, generation of synthetic
descriptions of the filtered set, and explanation of why two entities appear related through a
specific \textit{link}. This approach preserves user agency by keeping the engine's operations
visible and delegates only those tasks for which a language model adds real value.

\subsection*{Engine distribution across multiple nodes}

The engine currently runs as a single process against a single database. For larger-scale
collections, the system would benefit from decomposition into independent nodes: partitioning
storage by domain or \textit{dataset}, replicating facet computation to balance load, and
introducing a cache layer for frequent intermediate results. This line would require revisiting
the notion of session, today associated with a single process, and designing a coherent
invalidation protocol that maintains the immediacy on which the exploration experience rests.

\subsection*{Incremental \textit{query} generation}

In the current engine, each state contains by itself all the information needed to produce its
query, which is built in full at every interaction. An attractive alternative would be to
derive the new query from the previous one, reflecting at the technical level the cumulative
logic with which the user perceives the exploration. The motivation would not lie in
performance (the bulk of the work falls on executing the query against the database, not on its
construction, so any savings from avoiding regeneration would be marginal), but in having a
model and an implementation that are conceptually simpler and, therefore, more understandable.

That supposed simplicity, however, does not translate immediately to the technical level.
Moving to an incremental scheme would require introducing auxiliary support mechanisms (for
example, relying on database views or other forms of result reuse), each with its own
management problems when concurrent sessions come into play, as well as the need to guarantee
that the composition of queries remains correct in all cases. In the current design, each
query is determined solely by the state to which it belongs and is independent of the previous
ones, which already provides a simple model and implementation without assuming that added
complexity.

\subsection*{User layer, persistence, and shared history}

The introduction of an identity layer would allow state to be preserved between sessions,
previous explorations to be retrieved, union sets to be bookmarked as favorites, and
explorations to be shared through reproducible \textit{URLs}. This line is relatively
independent from the engine and could be developed in parallel, leveraging the existing
separation between back-end and front-end. Its greatest challenge is not technical, but rather
one of interaction design: integrating the new functionalities without overloading the current
interface, which has been tuned precisely to offer a light first experience focused on
exploration.

\subsection*{Extension to new domains and \textit{datasets}}

The system has been conceived from the start as domain-agnostic. The integration with TMDB and
DBLP has served to demonstrate that the relational model generalizes without substantial
modifications between an audiovisual catalog and a scientific publications database, but the
choice of future domains remains the most direct path to reinforce that evidence. Particularly
interesting would be complete library holdings (with various classes of objects coexisting
under the same descriptive schema), institutional document repositories (where the richness of
relations between versions, authors, and institutions would test the \textit{links} layer), and
open-source software catalogs (where dependency and authorship relations draw dense graphs). In
each case, effective incorporation requires adjusting the heuristics for facet identification
to the particularities of the domain and, in some cases, extending the loading process to
accommodate heterogeneous source formats.

Tied to the above, a natural improvement would be to incorporate into the \textit{API} an
\textbf{ingestion} flow exposed as such, replacing the external \textit{script} described in
appendix~\ref{Appendix:integracion}. Integrating the loading of new \textit{datasets} into the
service itself would simplify its use, open the door to automating the incorporation of new
holdings, and lower the barrier of entry for anyone wishing to explore domains different from
those validated.

\subsection*{Multiple value selection and disjunctive semantics}

In the current version, each facet admits the selection of a single value at a time, which
naturally leads to a \textbf{conjunctive} behavior (the result is the intersection of the
active filters over different facets). A reasonable extension consists of allowing
\textit{multiple selection} within the same facet, which would introduce a \textbf{disjunctive}
semantics between the chosen values while maintaining the conjunction between different facets.
The change is modest at the interface level, but it requires revisiting the state model and the
generation of \textit{queries} to accommodate the new operator, as well as reconsidering the
visual representation of active filters so that the distinction between \textit{and}/\textit{or}
is immediately legible.

\subsection*{Expansion of the unit test suite}

The validation carried out has relied fundamentally on integration tests over the TMDB and DBLP
dumps and on the qualitative usability evaluation. A clear incremental line that remains open
is the reinforcement of the \textbf{unit test suite}, especially over the engine components
that materialize the formal model's operations (state transitions, filter composition,
\textit{link} computation, and set operations). More systematic coverage at this level would
ease future refactorings and serve as a safety net against extensions of the model such as
those described in the previous subsections.

\subsection*{Incremental user experience improvements}

Several incremental improvements remain open on the developed interface, prioritized by the
usability evaluation: a version adapted to mobile devices, full internationalization beyond the
current Spanish-English pair, accessibility reinforcement through keyboard navigation and
\textit{ARIA} attributes, export of the navigation history and union set to reusable formats
(\textit{CSV}, \textit{JSON}, or bibliographic references), and more refined work on the visual
representation of the \textit{links} graph to make it useful with long paths.
